% =============================================================================
% Chapter 6 - Conclusion
% =============================================================================
% Included from main.tex via % =============================================================================
% Chapter 6 - Conclusion
% =============================================================================
% Included from main.tex via % =============================================================================
% Chapter 6 - Conclusion
% =============================================================================
% Included from main.tex via % =============================================================================
% Chapter 6 - Conclusion
% =============================================================================
% Included from main.tex via \include{contents/chapter6}.
%
% Rewritten 24 September 2026: opening framed on the aim and central research
% question of the task description; new Section 6.1 answering the research
% questions of Section 1.2; academic headings; no cross-references to other
% sections; citations only for claims about the reviewed literature, all from
% the verified bibliography. All figures quoted are those of Chapter 5.
% =============================================================================

\chapter{Conclusion}
\label{ch:conclusion}

This thesis investigated whether improved temporal modelling and motion representations can enhance the recognition of micro-expressions. Its central research question was how the modelling of temporal phases and motion cues can improve recognition accuracy for subtle and short-lived facial expressions. Building on established supervised models, the study evaluated practical strategies for temporal and motion modelling: a hybrid motion representation of optical flow and optical strain, Eulerian video magnification, a convolutional spatial stem, parameter-free attention and a temporal transformer based on self-attention. Rather than proposing a single architecture, it treated a standard recognition pipeline as a factorial experiment. Four components were switched on and off independently, giving twelve valid configurations, each trained from scratch and evaluated under complete 25-fold leave-one-subject-out cross-validation on 156 clips of CASME~II, with all other factors held identical.

This chapter answers the research questions, states the contributions of the thesis and the limitations within which they hold, outlines future work and closes with final remarks.


% =============================================================================
\section{Answers to the Research Questions}
\label{sec:rq-answers}
% =============================================================================

\textbf{Temporal modelling.} The first question asked to what extent temporal modelling with a transformer encoder improves recognition. It produced by far the largest improvement of the four components: pooled macro F1 rose by 0.2173 on average, and all six matched pairs improved. The best configuration, which contains the temporal transformer alone, reached a pooled macro F1 of 0.7122. The confusion matrices show that, without temporal modelling, the models over-predicted the minority classes and recognised few Negative clips, whereas with it they recognised all three classes more evenly. Because the transformer also adds learnable parameters, its effect cannot be attributed to temporal order alone.

\textbf{Motion magnification.} The second question asked whether amplifying subtle facial motion through Eulerian video magnification improves recognition. In this pipeline it did not do so reliably: its mean effect was +0.0152, four of six pairs improved, and adding it to the best configuration reduced the score by 0.0541. The clear gains reported for magnification with appearance-based descriptors~\cite{liX2018,bai2021} did not carry over to an input that already consists of optical flow and strain.

\textbf{Learned spatial features and attention.} The third question asked whether a learned convolutional stem and parameter-free attention add to the performance obtained from the hand-computed motion representation. Neither did: the stem had the only negative mean effect ($-0.0310$), including a loss of 0.1292 when added to the best configuration, and SimAM had a mean effect of $+0.0034$, indistinguishable from zero at this sample size.

\textbf{Central research question.} Taken together, the results indicate that, for a small dataset, the improvement comes from modelling the temporal evolution of an analytically computed motion representation, not from additional preprocessing or learned spatial processing. Modelling the temporal phases of the expression improved recognition, particularly of the minority classes, whereas the motion cues were best supplied directly by optical flow and strain. Robustness to noise was not measured separately, so the study cannot confirm that this improvement in sensitivity to fine-grained dynamics comes without a loss of robustness.


% =============================================================================
\section{Contributions}
\label{sec:contributions}
% =============================================================================

The thesis makes three contributions.

\textbf{Factorial evaluation of a standard pipeline.} The reviewed studies evaluate components within complete systems and vary them one at a time from the defaults of a single proposed design~\cite{liX2018,xia2020b,zhang2022}, so an improvement can rarely be attributed to the component credited with it. This thesis evaluated every architecturally valid combination of four components under one protocol, so that the contribution of each, including the learned components, could be measured against an otherwise identical pipeline.

\textbf{Distribution of effects across components.} Almost all of the measured effect is concentrated in one component, the temporal transformer. The convolutional stem reduces performance on average, SimAM cannot be distinguished from zero, and magnification has no consistent direction across its pairs. An end-to-end evaluation, which yields a single score, cannot reveal such a distribution.

\textbf{Performance of the complete pipeline.} The temporal transformer applied directly to the motion representation reached a pooled macro F1 of 0.7122, whereas the configuration with all four components reached 0.6659. An evaluation of the complete system alone would have reported the lower score, with no means of discovering that a subset of its own architecture performs better.


% =============================================================================
\section{Limitations}
\label{sec:ch6-limitations}
% =============================================================================

\textbf{Evaluation protocol.} Each configuration was trained once with a single random seed, so no confidence intervals are available for the differences between configurations. The training epoch was selected on the same held-out subject that was then scored, without an inner validation split, which biases the reported scores upwards by an unknown amount~\cite{varma2006}. Per-clip predictions were not retained, so no paired significance test between configurations can be computed~\cite{mcnemar1947}.

\textbf{Fixed preprocessing and training settings.} The magnification factor ($\alpha = 10$), the sequence length ($T = 32$), the SimAM constant ($\lambda = 10^{-4}$) and the focal loss exponent ($\gamma = 2$) were taken from the literature and held fixed, although such parameters have interior optima that can be located only by search~\cite{liX2018,ben2021}. No hyperparameter search was performed. Every effect is therefore conditional on these values, and each component was varied only as present or absent.

\textbf{Departures from the reviewed methods.} The implementation departs from the reviewed methods at twelve declared points. Three are most relevant to the results. The convolutional stem has no temporal kernels, so the ablation tests a learned spatial stem rather than spatio-temporal convolution. The stem operates at $224\times224$ pixels, above the resolutions favoured in the literature~\cite{xia2020b}, so its negative effect may reflect resolution as much as architecture. Magnification is applied after the clips are subsampled, so its realised pass-band depends on the clip length.

\textbf{Single dataset.} All results come from 156 clips of 25 subjects of one dataset, recorded in a laboratory from participants of one ethnicity and a narrow age range~\cite{yan2014,davison2018}. The models were trained from scratch, without pre-training or composite-dataset training. No claim therefore extends to other populations, recording conditions or datasets.


% =============================================================================
\section{Future Work}
\label{sec:further-work}
% =============================================================================

\textbf{Evaluation protocol.} Three changes would strengthen the evaluation without changing the model. Selecting the training epoch on subjects withheld from the training set, and scoring the held-out subject only once, would remove the selection bias of the reported scores. Retaining per-clip predictions would allow pairs of configurations to be compared with McNemar's test~\cite{mcnemar1947}. Repeating the experiment with several random seeds would provide confidence intervals for the effect of each component.

\textbf{Convolutional stem.} The most direct experiment is to retrain the configurations with the stem at lower resolutions, such as $28\times28$ and $56\times56$, within the range used by shallow micro-expression networks~\cite{liong2019b,xia2020b}. Varying kernel shape, input resolution and the position of the stem independently would separate the three explanations proposed for its weak effect: spatial-only kernels, an input that already encodes motion, and a high input resolution. Retaining the spatial layout of the stem's features, rather than averaging them over the whole frame, would also test whether the loss of location information is responsible.

\textbf{Motion representation.} Magnification should be applied before subsampling, so that the realised pass-band matches the intended 5--25 Hz, and the amplification factor should be varied for a motion-based input. Flow-based fine alignment~\cite{xu2017} and a comparison of the Farneb\"ack and TV-L1 estimators~\cite{zhaoS2021,liong2019a} remain untested. Adding a configuration without the strain channel would measure its contribution, and a dedicated strain filter and the four-term strain magnitude~\cite{liong2014a,liong2016} could be compared with the implemented form.

\textbf{Temporal modelling.} Three extensions follow from the published SLSTT architecture~\cite{zhang2022}: an LSTM aggregator in place of mean pooling, which its own ablation favours; flow computed relative to the onset rather than between consecutive frames; and learned rather than fixed positional embeddings~\cite{dosovitskiy2021}. The sequence length should be varied, and uniform sampling compared with the Temporal Interpolation Model~\cite{pfister2011}. Isolating self-attention from the added capacity of the transformer would require a control with a comparable number of parameters but no temporal modelling.

\textbf{Attention and training.} For SimAM, per-frame and per-clip statistics and different values of $\lambda$ should be compared~\cite{yang2021}. For training, focal loss could be compared with plain cross-entropy, and the focusing parameter and label smoothing varied. Because the configurations without the transformer over-predict the minority classes, adjusting the decision threshold to the true class priors after balanced sampling~\cite{buda2018} could also be evaluated.

\textbf{Interpretability and data.} Visual explanation of the trained models would show whether the transformer attends to the regions of the coded action units, and learning identity-invariant features would address the single-demographic limitation. Finally, evaluation on a composite dataset such as that of MEGC 2019~\cite{see2019}, and on larger or more diverse datasets such as SAMM, CAS(ME)$^3$ and DFME~\cite{davison2018,li2023casme3,zhao2024dfme}, is the direct test of whether the effects reported here hold beyond CASME~II.


% =============================================================================
\section{Final Remarks}
\label{sec:final-remarks}
% =============================================================================

Micro-expression recognition systems are assembled from stages that entered the field separately, each with its own reported gain, and the contribution of each stage within a complete pipeline has not been measured. A practitioner building such a system therefore has no evidence on which to decide where to invest effort. This thesis addressed that problem by evaluating every valid combination of four components under one protocol, so that each component could be assessed through matched pairs that differ in nothing else.

The results show that the stages do not contribute equally, and that the configuration that includes every component is not the one that performs best. These findings hold for one dataset, one training protocol and one setting of the preprocessing parameters, and they are stated with those conditions. Within them, they show that modelling the temporal evolution of facial motion matters more for micro-expression recognition than adding further preprocessing or spatial processing, and that a factorial evaluation can reveal what an evaluation of a single proposed system cannot.
.
%
% Rewritten 24 September 2026: opening framed on the aim and central research
% question of the task description; new Section 6.1 answering the research
% questions of Section 1.2; academic headings; no cross-references to other
% sections; citations only for claims about the reviewed literature, all from
% the verified bibliography. All figures quoted are those of Chapter 5.
% =============================================================================

\chapter{Conclusion}
\label{ch:conclusion}

This thesis investigated whether improved temporal modelling and motion representations can enhance the recognition of micro-expressions. Its central research question was how the modelling of temporal phases and motion cues can improve recognition accuracy for subtle and short-lived facial expressions. Building on established supervised models, the study evaluated practical strategies for temporal and motion modelling: a hybrid motion representation of optical flow and optical strain, Eulerian video magnification, a convolutional spatial stem, parameter-free attention and a temporal transformer based on self-attention. Rather than proposing a single architecture, it treated a standard recognition pipeline as a factorial experiment. Four components were switched on and off independently, giving twelve valid configurations, each trained from scratch and evaluated under complete 25-fold leave-one-subject-out cross-validation on 156 clips of CASME~II, with all other factors held identical.

This chapter answers the research questions, states the contributions of the thesis and the limitations within which they hold, outlines future work and closes with final remarks.


% =============================================================================
\section{Answers to the Research Questions}
\label{sec:rq-answers}
% =============================================================================

\textbf{Temporal modelling.} The first question asked to what extent temporal modelling with a transformer encoder improves recognition. It produced by far the largest improvement of the four components: pooled macro F1 rose by 0.2173 on average, and all six matched pairs improved. The best configuration, which contains the temporal transformer alone, reached a pooled macro F1 of 0.7122. The confusion matrices show that, without temporal modelling, the models over-predicted the minority classes and recognised few Negative clips, whereas with it they recognised all three classes more evenly. Because the transformer also adds learnable parameters, its effect cannot be attributed to temporal order alone.

\textbf{Motion magnification.} The second question asked whether amplifying subtle facial motion through Eulerian video magnification improves recognition. In this pipeline it did not do so reliably: its mean effect was +0.0152, four of six pairs improved, and adding it to the best configuration reduced the score by 0.0541. The clear gains reported for magnification with appearance-based descriptors~\cite{liX2018,bai2021} did not carry over to an input that already consists of optical flow and strain.

\textbf{Learned spatial features and attention.} The third question asked whether a learned convolutional stem and parameter-free attention add to the performance obtained from the hand-computed motion representation. Neither did: the stem had the only negative mean effect ($-0.0310$), including a loss of 0.1292 when added to the best configuration, and SimAM had a mean effect of $+0.0034$, indistinguishable from zero at this sample size.

\textbf{Central research question.} Taken together, the results indicate that, for a small dataset, the improvement comes from modelling the temporal evolution of an analytically computed motion representation, not from additional preprocessing or learned spatial processing. Modelling the temporal phases of the expression improved recognition, particularly of the minority classes, whereas the motion cues were best supplied directly by optical flow and strain. Robustness to noise was not measured separately, so the study cannot confirm that this improvement in sensitivity to fine-grained dynamics comes without a loss of robustness.


% =============================================================================
\section{Contributions}
\label{sec:contributions}
% =============================================================================

The thesis makes three contributions.

\textbf{Factorial evaluation of a standard pipeline.} The reviewed studies evaluate components within complete systems and vary them one at a time from the defaults of a single proposed design~\cite{liX2018,xia2020b,zhang2022}, so an improvement can rarely be attributed to the component credited with it. This thesis evaluated every architecturally valid combination of four components under one protocol, so that the contribution of each, including the learned components, could be measured against an otherwise identical pipeline.

\textbf{Distribution of effects across components.} Almost all of the measured effect is concentrated in one component, the temporal transformer. The convolutional stem reduces performance on average, SimAM cannot be distinguished from zero, and magnification has no consistent direction across its pairs. An end-to-end evaluation, which yields a single score, cannot reveal such a distribution.

\textbf{Performance of the complete pipeline.} The temporal transformer applied directly to the motion representation reached a pooled macro F1 of 0.7122, whereas the configuration with all four components reached 0.6659. An evaluation of the complete system alone would have reported the lower score, with no means of discovering that a subset of its own architecture performs better.


% =============================================================================
\section{Limitations}
\label{sec:ch6-limitations}
% =============================================================================

\textbf{Evaluation protocol.} Each configuration was trained once with a single random seed, so no confidence intervals are available for the differences between configurations. The training epoch was selected on the same held-out subject that was then scored, without an inner validation split, which biases the reported scores upwards by an unknown amount~\cite{varma2006}. Per-clip predictions were not retained, so no paired significance test between configurations can be computed~\cite{mcnemar1947}.

\textbf{Fixed preprocessing and training settings.} The magnification factor ($\alpha = 10$), the sequence length ($T = 32$), the SimAM constant ($\lambda = 10^{-4}$) and the focal loss exponent ($\gamma = 2$) were taken from the literature and held fixed, although such parameters have interior optima that can be located only by search~\cite{liX2018,ben2021}. No hyperparameter search was performed. Every effect is therefore conditional on these values, and each component was varied only as present or absent.

\textbf{Departures from the reviewed methods.} The implementation departs from the reviewed methods at twelve declared points. Three are most relevant to the results. The convolutional stem has no temporal kernels, so the ablation tests a learned spatial stem rather than spatio-temporal convolution. The stem operates at $224\times224$ pixels, above the resolutions favoured in the literature~\cite{xia2020b}, so its negative effect may reflect resolution as much as architecture. Magnification is applied after the clips are subsampled, so its realised pass-band depends on the clip length.

\textbf{Single dataset.} All results come from 156 clips of 25 subjects of one dataset, recorded in a laboratory from participants of one ethnicity and a narrow age range~\cite{yan2014,davison2018}. The models were trained from scratch, without pre-training or composite-dataset training. No claim therefore extends to other populations, recording conditions or datasets.


% =============================================================================
\section{Future Work}
\label{sec:further-work}
% =============================================================================

\textbf{Evaluation protocol.} Three changes would strengthen the evaluation without changing the model. Selecting the training epoch on subjects withheld from the training set, and scoring the held-out subject only once, would remove the selection bias of the reported scores. Retaining per-clip predictions would allow pairs of configurations to be compared with McNemar's test~\cite{mcnemar1947}. Repeating the experiment with several random seeds would provide confidence intervals for the effect of each component.

\textbf{Convolutional stem.} The most direct experiment is to retrain the configurations with the stem at lower resolutions, such as $28\times28$ and $56\times56$, within the range used by shallow micro-expression networks~\cite{liong2019b,xia2020b}. Varying kernel shape, input resolution and the position of the stem independently would separate the three explanations proposed for its weak effect: spatial-only kernels, an input that already encodes motion, and a high input resolution. Retaining the spatial layout of the stem's features, rather than averaging them over the whole frame, would also test whether the loss of location information is responsible.

\textbf{Motion representation.} Magnification should be applied before subsampling, so that the realised pass-band matches the intended 5--25 Hz, and the amplification factor should be varied for a motion-based input. Flow-based fine alignment~\cite{xu2017} and a comparison of the Farneb\"ack and TV-L1 estimators~\cite{zhaoS2021,liong2019a} remain untested. Adding a configuration without the strain channel would measure its contribution, and a dedicated strain filter and the four-term strain magnitude~\cite{liong2014a,liong2016} could be compared with the implemented form.

\textbf{Temporal modelling.} Three extensions follow from the published SLSTT architecture~\cite{zhang2022}: an LSTM aggregator in place of mean pooling, which its own ablation favours; flow computed relative to the onset rather than between consecutive frames; and learned rather than fixed positional embeddings~\cite{dosovitskiy2021}. The sequence length should be varied, and uniform sampling compared with the Temporal Interpolation Model~\cite{pfister2011}. Isolating self-attention from the added capacity of the transformer would require a control with a comparable number of parameters but no temporal modelling.

\textbf{Attention and training.} For SimAM, per-frame and per-clip statistics and different values of $\lambda$ should be compared~\cite{yang2021}. For training, focal loss could be compared with plain cross-entropy, and the focusing parameter and label smoothing varied. Because the configurations without the transformer over-predict the minority classes, adjusting the decision threshold to the true class priors after balanced sampling~\cite{buda2018} could also be evaluated.

\textbf{Interpretability and data.} Visual explanation of the trained models would show whether the transformer attends to the regions of the coded action units, and learning identity-invariant features would address the single-demographic limitation. Finally, evaluation on a composite dataset such as that of MEGC 2019~\cite{see2019}, and on larger or more diverse datasets such as SAMM, CAS(ME)$^3$ and DFME~\cite{davison2018,li2023casme3,zhao2024dfme}, is the direct test of whether the effects reported here hold beyond CASME~II.


% =============================================================================
\section{Final Remarks}
\label{sec:final-remarks}
% =============================================================================

Micro-expression recognition systems are assembled from stages that entered the field separately, each with its own reported gain, and the contribution of each stage within a complete pipeline has not been measured. A practitioner building such a system therefore has no evidence on which to decide where to invest effort. This thesis addressed that problem by evaluating every valid combination of four components under one protocol, so that each component could be assessed through matched pairs that differ in nothing else.

The results show that the stages do not contribute equally, and that the configuration that includes every component is not the one that performs best. These findings hold for one dataset, one training protocol and one setting of the preprocessing parameters, and they are stated with those conditions. Within them, they show that modelling the temporal evolution of facial motion matters more for micro-expression recognition than adding further preprocessing or spatial processing, and that a factorial evaluation can reveal what an evaluation of a single proposed system cannot.
.
%
% Rewritten 24 September 2026: opening framed on the aim and central research
% question of the task description; new Section 6.1 answering the research
% questions of Section 1.2; academic headings; no cross-references to other
% sections; citations only for claims about the reviewed literature, all from
% the verified bibliography. All figures quoted are those of Chapter 5.
% =============================================================================

\chapter{Conclusion}
\label{ch:conclusion}

This thesis investigated whether improved temporal modelling and motion representations can enhance the recognition of micro-expressions. Its central research question was how the modelling of temporal phases and motion cues can improve recognition accuracy for subtle and short-lived facial expressions. Building on established supervised models, the study evaluated practical strategies for temporal and motion modelling: a hybrid motion representation of optical flow and optical strain, Eulerian video magnification, a convolutional spatial stem, parameter-free attention and a temporal transformer based on self-attention. Rather than proposing a single architecture, it treated a standard recognition pipeline as a factorial experiment. Four components were switched on and off independently, giving twelve valid configurations, each trained from scratch and evaluated under complete 25-fold leave-one-subject-out cross-validation on 156 clips of CASME~II, with all other factors held identical.

This chapter answers the research questions, states the contributions of the thesis and the limitations within which they hold, outlines future work and closes with final remarks.


% =============================================================================
\section{Answers to the Research Questions}
\label{sec:rq-answers}
% =============================================================================

\textbf{Temporal modelling.} The first question asked to what extent temporal modelling with a transformer encoder improves recognition. It produced by far the largest improvement of the four components: pooled macro F1 rose by 0.2173 on average, and all six matched pairs improved. The best configuration, which contains the temporal transformer alone, reached a pooled macro F1 of 0.7122. The confusion matrices show that, without temporal modelling, the models over-predicted the minority classes and recognised few Negative clips, whereas with it they recognised all three classes more evenly. Because the transformer also adds learnable parameters, its effect cannot be attributed to temporal order alone.

\textbf{Motion magnification.} The second question asked whether amplifying subtle facial motion through Eulerian video magnification improves recognition. In this pipeline it did not do so reliably: its mean effect was +0.0152, four of six pairs improved, and adding it to the best configuration reduced the score by 0.0541. The clear gains reported for magnification with appearance-based descriptors~\cite{liX2018,bai2021} did not carry over to an input that already consists of optical flow and strain.

\textbf{Learned spatial features and attention.} The third question asked whether a learned convolutional stem and parameter-free attention add to the performance obtained from the hand-computed motion representation. Neither did: the stem had the only negative mean effect ($-0.0310$), including a loss of 0.1292 when added to the best configuration, and SimAM had a mean effect of $+0.0034$, indistinguishable from zero at this sample size.

\textbf{Central research question.} Taken together, the results indicate that, for a small dataset, the improvement comes from modelling the temporal evolution of an analytically computed motion representation, not from additional preprocessing or learned spatial processing. Modelling the temporal phases of the expression improved recognition, particularly of the minority classes, whereas the motion cues were best supplied directly by optical flow and strain. Robustness to noise was not measured separately, so the study cannot confirm that this improvement in sensitivity to fine-grained dynamics comes without a loss of robustness.


% =============================================================================
\section{Contributions}
\label{sec:contributions}
% =============================================================================

The thesis makes three contributions.

\textbf{Factorial evaluation of a standard pipeline.} The reviewed studies evaluate components within complete systems and vary them one at a time from the defaults of a single proposed design~\cite{liX2018,xia2020b,zhang2022}, so an improvement can rarely be attributed to the component credited with it. This thesis evaluated every architecturally valid combination of four components under one protocol, so that the contribution of each, including the learned components, could be measured against an otherwise identical pipeline.

\textbf{Distribution of effects across components.} Almost all of the measured effect is concentrated in one component, the temporal transformer. The convolutional stem reduces performance on average, SimAM cannot be distinguished from zero, and magnification has no consistent direction across its pairs. An end-to-end evaluation, which yields a single score, cannot reveal such a distribution.

\textbf{Performance of the complete pipeline.} The temporal transformer applied directly to the motion representation reached a pooled macro F1 of 0.7122, whereas the configuration with all four components reached 0.6659. An evaluation of the complete system alone would have reported the lower score, with no means of discovering that a subset of its own architecture performs better.


% =============================================================================
\section{Limitations}
\label{sec:ch6-limitations}
% =============================================================================

\textbf{Evaluation protocol.} Each configuration was trained once with a single random seed, so no confidence intervals are available for the differences between configurations. The training epoch was selected on the same held-out subject that was then scored, without an inner validation split, which biases the reported scores upwards by an unknown amount~\cite{varma2006}. Per-clip predictions were not retained, so no paired significance test between configurations can be computed~\cite{mcnemar1947}.

\textbf{Fixed preprocessing and training settings.} The magnification factor ($\alpha = 10$), the sequence length ($T = 32$), the SimAM constant ($\lambda = 10^{-4}$) and the focal loss exponent ($\gamma = 2$) were taken from the literature and held fixed, although such parameters have interior optima that can be located only by search~\cite{liX2018,ben2021}. No hyperparameter search was performed. Every effect is therefore conditional on these values, and each component was varied only as present or absent.

\textbf{Departures from the reviewed methods.} The implementation departs from the reviewed methods at twelve declared points. Three are most relevant to the results. The convolutional stem has no temporal kernels, so the ablation tests a learned spatial stem rather than spatio-temporal convolution. The stem operates at $224\times224$ pixels, above the resolutions favoured in the literature~\cite{xia2020b}, so its negative effect may reflect resolution as much as architecture. Magnification is applied after the clips are subsampled, so its realised pass-band depends on the clip length.

\textbf{Single dataset.} All results come from 156 clips of 25 subjects of one dataset, recorded in a laboratory from participants of one ethnicity and a narrow age range~\cite{yan2014,davison2018}. The models were trained from scratch, without pre-training or composite-dataset training. No claim therefore extends to other populations, recording conditions or datasets.


% =============================================================================
\section{Future Work}
\label{sec:further-work}
% =============================================================================

\textbf{Evaluation protocol.} Three changes would strengthen the evaluation without changing the model. Selecting the training epoch on subjects withheld from the training set, and scoring the held-out subject only once, would remove the selection bias of the reported scores. Retaining per-clip predictions would allow pairs of configurations to be compared with McNemar's test~\cite{mcnemar1947}. Repeating the experiment with several random seeds would provide confidence intervals for the effect of each component.

\textbf{Convolutional stem.} The most direct experiment is to retrain the configurations with the stem at lower resolutions, such as $28\times28$ and $56\times56$, within the range used by shallow micro-expression networks~\cite{liong2019b,xia2020b}. Varying kernel shape, input resolution and the position of the stem independently would separate the three explanations proposed for its weak effect: spatial-only kernels, an input that already encodes motion, and a high input resolution. Retaining the spatial layout of the stem's features, rather than averaging them over the whole frame, would also test whether the loss of location information is responsible.

\textbf{Motion representation.} Magnification should be applied before subsampling, so that the realised pass-band matches the intended 5--25 Hz, and the amplification factor should be varied for a motion-based input. Flow-based fine alignment~\cite{xu2017} and a comparison of the Farneb\"ack and TV-L1 estimators~\cite{zhaoS2021,liong2019a} remain untested. Adding a configuration without the strain channel would measure its contribution, and a dedicated strain filter and the four-term strain magnitude~\cite{liong2014a,liong2016} could be compared with the implemented form.

\textbf{Temporal modelling.} Three extensions follow from the published SLSTT architecture~\cite{zhang2022}: an LSTM aggregator in place of mean pooling, which its own ablation favours; flow computed relative to the onset rather than between consecutive frames; and learned rather than fixed positional embeddings~\cite{dosovitskiy2021}. The sequence length should be varied, and uniform sampling compared with the Temporal Interpolation Model~\cite{pfister2011}. Isolating self-attention from the added capacity of the transformer would require a control with a comparable number of parameters but no temporal modelling.

\textbf{Attention and training.} For SimAM, per-frame and per-clip statistics and different values of $\lambda$ should be compared~\cite{yang2021}. For training, focal loss could be compared with plain cross-entropy, and the focusing parameter and label smoothing varied. Because the configurations without the transformer over-predict the minority classes, adjusting the decision threshold to the true class priors after balanced sampling~\cite{buda2018} could also be evaluated.

\textbf{Interpretability and data.} Visual explanation of the trained models would show whether the transformer attends to the regions of the coded action units, and learning identity-invariant features would address the single-demographic limitation. Finally, evaluation on a composite dataset such as that of MEGC 2019~\cite{see2019}, and on larger or more diverse datasets such as SAMM, CAS(ME)$^3$ and DFME~\cite{davison2018,li2023casme3,zhao2024dfme}, is the direct test of whether the effects reported here hold beyond CASME~II.


% =============================================================================
\section{Final Remarks}
\label{sec:final-remarks}
% =============================================================================

Micro-expression recognition systems are assembled from stages that entered the field separately, each with its own reported gain, and the contribution of each stage within a complete pipeline has not been measured. A practitioner building such a system therefore has no evidence on which to decide where to invest effort. This thesis addressed that problem by evaluating every valid combination of four components under one protocol, so that each component could be assessed through matched pairs that differ in nothing else.

The results show that the stages do not contribute equally, and that the configuration that includes every component is not the one that performs best. These findings hold for one dataset, one training protocol and one setting of the preprocessing parameters, and they are stated with those conditions. Within them, they show that modelling the temporal evolution of facial motion matters more for micro-expression recognition than adding further preprocessing or spatial processing, and that a factorial evaluation can reveal what an evaluation of a single proposed system cannot.
.
%
% Rewritten 24 September 2026: opening framed on the aim and central research
% question of the task description; new Section 6.1 answering the research
% questions of Section 1.2; academic headings; no cross-references to other
% sections; citations only for claims about the reviewed literature, all from
% the verified bibliography. All figures quoted are those of Chapter 5.
% =============================================================================

\chapter{Conclusion}
\label{ch:conclusion}

This thesis investigated whether improved temporal modelling and motion representations can enhance the recognition of micro-expressions. Its central research question was how the modelling of temporal phases and motion cues can improve recognition accuracy for subtle and short-lived facial expressions. Building on established supervised models, the study evaluated practical strategies for temporal and motion modelling: a hybrid motion representation of optical flow and optical strain, Eulerian video magnification, a convolutional spatial stem, parameter-free attention and a temporal transformer based on self-attention. Rather than proposing a single architecture, it treated a standard recognition pipeline as a factorial experiment. Four components were switched on and off independently, giving twelve valid configurations, each trained from scratch and evaluated under complete 25-fold leave-one-subject-out cross-validation on 156 clips of CASME~II, with all other factors held identical.

This chapter answers the research questions, states the contributions of the thesis and the limitations within which they hold, outlines future work and closes with final remarks.


% =============================================================================
\section{Answers to the Research Questions}
\label{sec:rq-answers}
% =============================================================================

\textbf{Temporal modelling.} The first question asked to what extent temporal modelling with a transformer encoder improves recognition. It produced by far the largest improvement of the four components: pooled macro F1 rose by 0.2173 on average, and all six matched pairs improved. The best configuration, which contains the temporal transformer alone, reached a pooled macro F1 of 0.7122. The confusion matrices show that, without temporal modelling, the models over-predicted the minority classes and recognised few Negative clips, whereas with it they recognised all three classes more evenly. Because the transformer also adds learnable parameters, its effect cannot be attributed to temporal order alone.

\textbf{Motion magnification.} The second question asked whether amplifying subtle facial motion through Eulerian video magnification improves recognition. In this pipeline it did not do so reliably: its mean effect was +0.0152, four of six pairs improved, and adding it to the best configuration reduced the score by 0.0541. The clear gains reported for magnification with appearance-based descriptors~\cite{liX2018,bai2021} did not carry over to an input that already consists of optical flow and strain.

\textbf{Learned spatial features and attention.} The third question asked whether a learned convolutional stem and parameter-free attention add to the performance obtained from the hand-computed motion representation. Neither did: the stem had the only negative mean effect ($-0.0310$), including a loss of 0.1292 when added to the best configuration, and SimAM had a mean effect of $+0.0034$, indistinguishable from zero at this sample size.

\textbf{Central research question.} Taken together, the results indicate that, for a small dataset, the improvement comes from modelling the temporal evolution of an analytically computed motion representation, not from additional preprocessing or learned spatial processing. Modelling the temporal phases of the expression improved recognition, particularly of the minority classes, whereas the motion cues were best supplied directly by optical flow and strain. Robustness to noise was not measured separately, so the study cannot confirm that this improvement in sensitivity to fine-grained dynamics comes without a loss of robustness.


% =============================================================================
\section{Contributions}
\label{sec:contributions}
% =============================================================================

The thesis makes three contributions.

\textbf{Factorial evaluation of a standard pipeline.} The reviewed studies evaluate components within complete systems and vary them one at a time from the defaults of a single proposed design~\cite{liX2018,xia2020b,zhang2022}, so an improvement can rarely be attributed to the component credited with it. This thesis evaluated every architecturally valid combination of four components under one protocol, so that the contribution of each, including the learned components, could be measured against an otherwise identical pipeline.

\textbf{Distribution of effects across components.} Almost all of the measured effect is concentrated in one component, the temporal transformer. The convolutional stem reduces performance on average, SimAM cannot be distinguished from zero, and magnification has no consistent direction across its pairs. An end-to-end evaluation, which yields a single score, cannot reveal such a distribution.

\textbf{Performance of the complete pipeline.} The temporal transformer applied directly to the motion representation reached a pooled macro F1 of 0.7122, whereas the configuration with all four components reached 0.6659. An evaluation of the complete system alone would have reported the lower score, with no means of discovering that a subset of its own architecture performs better.


% =============================================================================
\section{Limitations}
\label{sec:ch6-limitations}
% =============================================================================

\textbf{Evaluation protocol.} Each configuration was trained once with a single random seed, so no confidence intervals are available for the differences between configurations. The training epoch was selected on the same held-out subject that was then scored, without an inner validation split, which biases the reported scores upwards by an unknown amount~\cite{varma2006}. Per-clip predictions were not retained, so no paired significance test between configurations can be computed~\cite{mcnemar1947}.

\textbf{Fixed preprocessing and training settings.} The magnification factor ($\alpha = 10$), the sequence length ($T = 32$), the SimAM constant ($\lambda = 10^{-4}$) and the focal loss exponent ($\gamma = 2$) were taken from the literature and held fixed, although such parameters have interior optima that can be located only by search~\cite{liX2018,ben2021}. No hyperparameter search was performed. Every effect is therefore conditional on these values, and each component was varied only as present or absent.

\textbf{Departures from the reviewed methods.} The implementation departs from the reviewed methods at twelve declared points. Three are most relevant to the results. The convolutional stem has no temporal kernels, so the ablation tests a learned spatial stem rather than spatio-temporal convolution. The stem operates at $224\times224$ pixels, above the resolutions favoured in the literature~\cite{xia2020b}, so its negative effect may reflect resolution as much as architecture. Magnification is applied after the clips are subsampled, so its realised pass-band depends on the clip length.

\textbf{Single dataset.} All results come from 156 clips of 25 subjects of one dataset, recorded in a laboratory from participants of one ethnicity and a narrow age range~\cite{yan2014,davison2018}. The models were trained from scratch, without pre-training or composite-dataset training. No claim therefore extends to other populations, recording conditions or datasets.


% =============================================================================
\section{Future Work}
\label{sec:further-work}
% =============================================================================

\textbf{Evaluation protocol.} Three changes would strengthen the evaluation without changing the model. Selecting the training epoch on subjects withheld from the training set, and scoring the held-out subject only once, would remove the selection bias of the reported scores. Retaining per-clip predictions would allow pairs of configurations to be compared with McNemar's test~\cite{mcnemar1947}. Repeating the experiment with several random seeds would provide confidence intervals for the effect of each component.

\textbf{Convolutional stem.} The most direct experiment is to retrain the configurations with the stem at lower resolutions, such as $28\times28$ and $56\times56$, within the range used by shallow micro-expression networks~\cite{liong2019b,xia2020b}. Varying kernel shape, input resolution and the position of the stem independently would separate the three explanations proposed for its weak effect: spatial-only kernels, an input that already encodes motion, and a high input resolution. Retaining the spatial layout of the stem's features, rather than averaging them over the whole frame, would also test whether the loss of location information is responsible.

\textbf{Motion representation.} Magnification should be applied before subsampling, so that the realised pass-band matches the intended 5--25 Hz, and the amplification factor should be varied for a motion-based input. Flow-based fine alignment~\cite{xu2017} and a comparison of the Farneb\"ack and TV-L1 estimators~\cite{zhaoS2021,liong2019a} remain untested. Adding a configuration without the strain channel would measure its contribution, and a dedicated strain filter and the four-term strain magnitude~\cite{liong2014a,liong2016} could be compared with the implemented form.

\textbf{Temporal modelling.} Three extensions follow from the published SLSTT architecture~\cite{zhang2022}: an LSTM aggregator in place of mean pooling, which its own ablation favours; flow computed relative to the onset rather than between consecutive frames; and learned rather than fixed positional embeddings~\cite{dosovitskiy2021}. The sequence length should be varied, and uniform sampling compared with the Temporal Interpolation Model~\cite{pfister2011}. Isolating self-attention from the added capacity of the transformer would require a control with a comparable number of parameters but no temporal modelling.

\textbf{Attention and training.} For SimAM, per-frame and per-clip statistics and different values of $\lambda$ should be compared~\cite{yang2021}. For training, focal loss could be compared with plain cross-entropy, and the focusing parameter and label smoothing varied. Because the configurations without the transformer over-predict the minority classes, adjusting the decision threshold to the true class priors after balanced sampling~\cite{buda2018} could also be evaluated.

\textbf{Interpretability and data.} Visual explanation of the trained models would show whether the transformer attends to the regions of the coded action units, and learning identity-invariant features would address the single-demographic limitation. Finally, evaluation on a composite dataset such as that of MEGC 2019~\cite{see2019}, and on larger or more diverse datasets such as SAMM, CAS(ME)$^3$ and DFME~\cite{davison2018,li2023casme3,zhao2024dfme}, is the direct test of whether the effects reported here hold beyond CASME~II.


% =============================================================================
\section{Final Remarks}
\label{sec:final-remarks}
% =============================================================================

Micro-expression recognition systems are assembled from stages that entered the field separately, each with its own reported gain, and the contribution of each stage within a complete pipeline has not been measured. A practitioner building such a system therefore has no evidence on which to decide where to invest effort. This thesis addressed that problem by evaluating every valid combination of four components under one protocol, so that each component could be assessed through matched pairs that differ in nothing else.

The results show that the stages do not contribute equally, and that the configuration that includes every component is not the one that performs best. These findings hold for one dataset, one training protocol and one setting of the preprocessing parameters, and they are stated with those conditions. Within them, they show that modelling the temporal evolution of facial motion matters more for micro-expression recognition than adding further preprocessing or spatial processing, and that a factorial evaluation can reveal what an evaluation of a single proposed system cannot.
